\documentclass[11pt,a4paper]{article}
\usepackage[utf8]{inputenc}
\usepackage[T1]{fontenc}
\usepackage[english]{babel}
\usepackage{amsmath,amssymb}
\usepackage{geometry}
\usepackage{hyperref}
\geometry{margin=2.2cm}

\title{ADPro for Robotic Diffusion Policies\\\large Algorithm Steps and Code-Level Implementation}
\author{SIMUL Project -- Franka Panda}
\date{\today}

\begin{document}
\maketitle

\section{Objective}
ADPro (\emph{Adaptive Diffusion Policy}) is a \emph{test-time} adaptation method for a pre-trained diffusion policy.
The main idea is to avoid re-training and instead guide denoising using task geometry.

In this project, ADPro is implemented in \texttt{models/adpro.py} through the \texttt{ADPro} class, with two variants:
\begin{itemize}
\item \texttt{implementation='paper'}: most faithful version to the paper (FGR + Chamfer guidance + spherical constraint).
\item \texttt{implementation='practical'}: more robust/pragmatic simulation version (explicit target + optional blend).
\end{itemize}

\section{Notation}
\begin{itemize}
\item Observation: $O = (O^0, O^1)$ where $O^0$ is the gripper point cloud and $O^1$ is the scene point cloud.
\item 7D action: $a=[x,y,z,q_x,q_y,q_z,q_w]$.
\item Diffusion step: $t$ (from high noise to low noise).
\item Base policy: network $\epsilon_\theta(O,a_t,t)$.
\end{itemize}

\section{Step 1 -- Task-Aware Initialization}
\subsection{Theory}
Instead of starting from pure Gaussian noise $a_T\sim\mathcal N(0,I)$, ADPro computes a structured initial pose $a_0^{\text{est}}$ from scene geometry,
then projects it to noise level $t=M$:
\begin{equation}
 a_M = \sqrt{\bar\alpha_M}\,a_0^{\text{est}} + \sqrt{1-\bar\alpha_M}\,\epsilon,\quad \epsilon\sim\mathcal N(0,I).
\end{equation}

\subsection{Implementation}
Function: \texttt{\_task\_aware\_init(obs, t\_start)}.
\begin{itemize}
\item Point subsampling: \texttt{subsample\_pointcloud}.
\item Object-point filtering (to reduce table influence): $z$-based mask.
\item \texttt{paper} variant: full rigid estimate with \texttt{fast\_global\_registration}, then \texttt{se3\_to\_action}.
\item \texttt{practical} variant: explicit grasp target (via \texttt{\_target\_grasp\_action}) + centroid-based correction.
\end{itemize}

\section{Step 2 -- One Diffusion Denoising Step}
\subsection{Principle}
The standard DDPM reverse step is first applied to obtain $\tilde a_{t-1}$.

\subsection{Implementation}
In \texttt{\_adpro\_step}:
\begin{itemize}
\item Noise prediction: \texttt{eps\_pred = policy.network.forward(obs, a\_t, t)}.
\item DDPM step: \texttt{a\_tilde = policy.ddpm\_step(...)}.
\item Deterministic option: \texttt{deterministic\_denoising=True} removes reverse-process stochastic noise for stability.
\end{itemize}

\section{Step 3 -- Clean Action Estimation (Tweedie-like Formula)}
\begin{equation}
\hat a_0 = \frac{a_t - \sqrt{1-\bar\alpha_t}\,\epsilon_\theta(O,a_t,t)}{\sqrt{\bar\alpha_t}}.
\end{equation}
Implementation: \texttt{\_predict\_clean\_action}.

\section{Step 4 -- Task-Manifold Guidance}
\subsection{Paper Version (surface-to-surface)}
The geometric objective is built on point clouds (Chamfer):
\begin{equation}
\mathcal L = d_{\text{Chamfer}}\big(T(\hat a_0)\,O^0,\;O^1\big).
\end{equation}
The intermediate action is then corrected in the opposite gradient direction.

Implementation in \texttt{\_apply\_task\_manifold\_guidance}:
\begin{itemize}
\item Action $\to$ pose conversion: \texttt{action\_to\_se3}.
\item Back-transform to local gripper frame, then apply $T(\hat a_0)$ (frame consistency).
\item Chamfer gradient: \texttt{chamfer\_distance\_grad}.
\end{itemize}

\subsection{Practical Version (point-to-point)}
Guidance pushes toward an explicit \texttt{target\_grasp\_action} (target position + quaternion) using a simplified L2 error.

\section{Step 5 -- Spherical Constraint (Eq. 9)}
The correction is normalized and scaled by the current noise level:
\begin{equation}
 a_{t-1} = \tilde a_{t-1} - s\,\sqrt{d}\,\sigma_t\,\frac{\nabla_a\mathcal L}{\|\nabla_a\mathcal L\|},
\end{equation}
where $d$ is the action dimension and $s$ is a stabilization factor (in code: \texttt{spherical\_scale}).

Implementation:
\begin{itemize}
\item \texttt{use\_spherical=True}: enables the normalized update rule.
\item \texttt{spherical\_scale}: controls guidance step intensity.
\end{itemize}

\section{Step 6 -- Full Inference Loop}
In \texttt{inference(...)}:
\begin{enumerate}
\item Choose initialization: task-aware (if \texttt{use\_fgr}) or Gaussian.
\item Iterate decreasing $t$ from \texttt{t\_start} down to $0$.
\item At each step: DDPM $\to$ clean estimate $\hat a_0$ $\to$ manifold guidance $\to$ spherical constraint.
\item Final output: action $a_0$ applied to the environment/controller.
\end{enumerate}

\section{Main Parameters and Their Role}
\begin{itemize}
\item \texttt{implementation}: \texttt{paper} vs \texttt{practical}.
\item \texttt{M}: initial noise level (smaller means closer-to-final initialization).
\item \texttt{spherical\_scale}: stability vs aggressiveness of guidance.
\item \texttt{deterministic\_denoising}: temporal stability in simulation and robot loops.
\item \texttt{use\_blend} (practical): additional stabilization toward explicit target.
\end{itemize}

\section{Visualization and Evaluation in This Repository}
\subsection{Simulation and Offline Evaluation}
In \texttt{simulate\_realtime.py}, \texttt{--show-denoise} visualizes:
\begin{itemize}
\item internal denoising trajectory $a_t\to a_0$ (baseline vs ADPro),
\item ADPro initialization state,
\item guidance metrics (gradient norm, step norm, blend).
\end{itemize}

In \texttt{demo.py}, ADPro is benchmarked against baseline with:
\begin{itemize}
\item denoising convergence curves,
\item action-component/backtracking plots,
\item 3D executed trajectories,
\item MSE vs denoising steps,
\item success-rate vs denoising steps.
\end{itemize}

\subsection{Real-Robot Layer (Updated)}
The repository now includes a separated deployment layer under \texttt{real\_robot/}:
\begin{itemize}
\item \texttt{control\_loop.py}, \texttt{policy\_engine.py}, \texttt{safety.py}: observation-policy-command loop with safety guards.
\item Adapters: \texttt{mock\_adapter.py} (local dry-run), \texttt{ros2\_franka\_adapter.py} (ROS2), \texttt{gello\_zmq\_adapter.py} (GELLO/ZMQ stack).
\item Operational scripts: \texttt{run\_policy.py}, \texttt{calibrate\_workspace\_guided.py}, \texttt{calibrate\_local\_frame\_matrix.py}, and motion-test scripts.
\item Config files: \texttt{real\_robot/config/lab\_franka.json} and \texttt{lab\_franka.yaml}.
\end{itemize}

\subsection{Calibration and Deployment Notes}
Calibration is required to align camera/local coordinates with robot/world coordinates before deploying ADPro on hardware.
The added scripts support workspace-bound estimation and local-frame matrix calibration to make execution reproducible and safe.

\section{Summary}
ADPro improves a diffusion policy by injecting geometric structure at test time:
task-aware initialization, manifold guidance, and spherical-normalized updates.
The \texttt{paper} variant maximizes theoretical fidelity, while the \texttt{practical} variant maximizes robustness in simulation.

\end{document}
